\documentclass[11pt,a4paper]{article}
\usepackage[utf8]{inputenc}
\usepackage[spanish]{babel}
\usepackage{amsmath}
\usepackage{amsfonts}
\usepackage{booktabs}
\usepackage{graphicx}
\usepackage{listings}
\usepackage{xcolor}
\usepackage{geometry}
\geometry{margin=1in}

% Configuración para bloques de código
\definecolor{mGreen}{rgb}{0,0.6,0}
\definecolor{mGray}{rgb}{0.5,0.5,0.5}
\definecolor{mPurple}{rgb}{0.58,0,0.82}
\definecolor{backgroundColour}{rgb}{0.95,0.95,0.92}

\lstset{
    backgroundcolor=\color{backgroundColour},
    commentstyle=\color{mGreen},
    keywordstyle=\color{magenta},
    numberstyle=\tiny\color{mGray},
    stringstyle=\color{mPurple},
    basicstyle=\ttfamily\footnotesize,
    breakatwhitespace=false,
    breaklines=true,
    captionpos=b,
    keepspaces=true,
    numbers=left,
    numbersep=5pt,
    showspaces=false,
    showstringspaces=false,
    showtabs=false,
    tabsize=2,
    language=C++
}

\title{Informe de Programación Paralela: Balanceo de Carga y Sincronización}
\author{Ciencias de la Computación}
\date{21 de mayo de 2026}

\begin{document}

\maketitle

\section{Ejercicio 3: Aproximación de la Función Zeta de Riemann}
El objetivo de este ejercicio fue analizar e implementar la paralelización de un bucle con complejidad temporal $O(k^2)$ mediante el uso de hilos estándar de C++11 (\texttt{std::thread}).

\subsection{Análisis de Dependencias y Variables Compartidas}
\begin{itemize}
    \item \textbf{Dependencias de Datos:} No existen dependencias de datos entre las iteraciones del bucle externo (\textit{loop-carried dependencies}). Cada elemento $X[k]$ se calcula de forma independiente, clasificando el problema como \textit{Embarrassingly Parallel}.
    \item \textbf{Variables Compartidas:} El vector de resultados $X$ es compartido por todos los hilos. No obstante, dado que cada hilo escribe en índices estrictamente disjuntos determinados por el patrón de distribución, no se generan condiciones de carrera (\textit{data races}), haciendo innecesario el uso de mecanismos de sincronización costosos.
\end{itemize}

\subsection{Patrones de Distribución de Carga}
Se evaluaron tres patrones de diseño fundamentales para la asignación de trabajo:
\begin{enumerate}
    \item \textbf{Parallel Block (Bloques):} Divide el dominio $N$ en subrangos continuos uniformes.
    \item \textbf{Parallel Cyclic (Cíclico):} Distribuye las iteraciones de manera intercalada (Round Robin) de uno en uno.
    \item \textbf{Parallel Block Cyclic (Bloque-Cíclico):} Variante híbrida que distribuye bloques de tamaño fijo ($BS$) de forma cíclica.
\end{enumerate}

\subsection{Evaluación Experimental y Resultados}
La asimetría del cálculo (las iteraciones con $k$ alto requieren drásticamente más tiempo que las iniciales) genera un severo desbalanceo de carga en la distribución por bloques. A continuación se presentan los tiempos medidos en milisegundos:

\begin{table}[h!]
\centering
\caption{Resultados de las Pruebas de Rendimiento (en ms)}
\begin{tabular}{lccccc}
\toprule
\textbf{Configuración} & \textbf{Hilos} & \textbf{Size $N$} & \textbf{Cíclico} & \textbf{Bloques} & \textbf{Bloque-Cíclico} \\ \midrule
Prueba 1 & 4 & 2000 & 12,227 & 24,442 & 14,603 (BS=50) \\
Prueba 2 & 8 & 4000 & 56,291 & 134,588 & 71,269 (BS=100) \\
\bottomrule
\end{tabular}
\end{table}

\begin{table}[h!]
\centering
\caption{Prueba 3: Impacto del Tamaño de Bloque en Bloque-Cíclico ($N=3500$, 8 Hilos)}
\begin{tabular}{lcccc}
\toprule
\textbf{Métrica / Variación} & \textbf{Cíclico Puro} & \textbf{Bloques Puros} & \textbf{Bloque-Cíclico} & \textbf{Tamaño Bloque (BS)} \\ \midrule
Prueba 3.1 & 39,769 ms & 87,735 ms & 41,360 ms & BS = 1 \\
Prueba 3.2 & 39,928 ms & 88,603 ms & 84,923 ms & BS = 437 \\
Prueba 3.3 & 39,671 ms & 85,811 ms & 43,027 ms & BS = 50 \\
\bottomrule
\end{tabular}
\end{table}

\textbf{Conclusión del experimento:} El patrón Cíclico minimiza los tiempos al garantizar que el "sufrimiento computacional" de los valores de $k$ más altos se reparta equitativamente. La prueba 3.1 demuestra que con $BS=1$ el comportamiento es idéntico al Cíclico, mientras que la prueba 3.2 con $BS = N/Hilos$ degenera en el patrón por Bloques. Al tratarse de un problema limitado por procesamiento (\textit{CPU-bound}), el balanceo puro supera a la optimización de memoria caché.

\newpage
\section{Ejercicio 4: Sincronización con Variables de Condición}
Se diseñó un mecanismo de unión (\textit{custom join}) manual para detener el hilo principal hasta la completa finalización del hilo hijo utilizando variables de condición.

\subsection{Implementación de Bloqueos (Scoped vs. Explicit)}
Se contrastaron dos enfoques de exclusión mutua para la actualización de la variable de estado \texttt{done}:
\begin{itemize}
    \item \textbf{Explicit Locking (\texttt{m.lock() / m.unlock()}):} Requiere la liberación manual del recurso. Presenta vulnerabilidades ante posibles excepciones en el flujo que impidan llegar al \texttt{unlock()}, provocando interbloqueos (\textit{deadlocks}).
    \item \textbf{Scoped Locking (\texttt{std::lock\_guard}):} Utiliza el patrón RAII. Garantiza la liberación segura y automática del mutex al salir del bloque de alcance (\textit{scope}) delimitado por llaves.
\end{itemize}

Ambas implementaciones arrojaron el orden de salida esperado de manera consistente:
\begin{lstlisting}[language=bash]
child
parent
\end{lstlisting}

\section{Ejercicio 5: Sincronización One-Shot usando Futures y Promises}
Para escenarios donde la comunicación ocurre una única vez (sincronización por ocurrencia de evento único o \textit{One-Shot}), las variables de condición introducen una complejidad innecesaria.

A través de la sección 4.5 del texto guía, se reestructuró el programa utilizando \texttt{std::promise} y \texttt{std::future}. Este patrón elimina por completo las variables globales de control, los mutexes explícitos y la lógica de bucles condicionales del lado del receptor. El hilo secundario invoca \texttt{p.set\_value()} (disparo único) desbloqueando limpiamente la barrera de espera del hilo principal configurada mediante \texttt{f.wait()}.

\end{document}
